\documentclass{article}

% Language setting
% Replace `english' with e.g. `spanish' to change the document language
\usepackage[english]{babel}
\usepackage{amssymb}
% Set page size and margins
% Replace `letterpaper' with `a4paper' for UK/EU standard size
\usepackage[letterpaper,top=2cm,bottom=2cm,left=3cm,right=3cm,marginparwidth=1.75cm]{geometry}

% Useful packages
\usepackage{amsmath}
\usepackage{amsthm}
\usepackage{graphicx}
\usepackage[colorlinks=true, allcolors=blue]{hyperref}


%%%%%%%%%%%%%%%%%%%%%%%%%%%%%%%%%%%%%%%%%%%%%%%%%%%%%%%%%%%%%%%%%%%%

\newtheorem{example}{Example} 
\newtheorem{theorem}{Theorem}
\newtheorem{lemma}[theorem]{Lemma} 
\newtheorem{Theorem}[theorem]{Theorem} 
\newtheorem{remark}[theorem]{Remark}
\newtheorem{corollary}[theorem]{Corollary}
\newtheorem{definition}[theorem]{Definition}
\newtheorem{conjecture}[theorem]{Conjecture}
\newtheorem{axiom}[theorem]{Axiom}
\newtheorem{proposition}{Proposition}
\newtheorem{condition}{Condition}
\newenvironment{proofof}[1]{\begin{proof}\textbf{of {#1}}}{\end{proof}}

%%%%%%%%%%%%%%%%%%%%%%%%%%%%%%%%%%%%%%%%%%%%%%%%%%%%%%%%%%%%%%%%%%%%



\title{Easton+: Embedding-based Easton algorithm}
\author{You}

\begin{document}
\maketitle

\section{Introduction}

In the context of Retrieval-Augmented Generation (RAG), accurately attributing a Large Language Model's (LLM) generated response to specific retrieved snippets is a critical challenge for transparency and maintaining a healthy WEB ecosystem. The Easton algorithm addresses this by leveraging Shapley values to distribute attribution scores across retrieved snippets based on their marginal contribution to the LLM's response log-likelihood. However, relying directly on LLM's likelihoods introduces severe limitations: the approach is computationally intractable requiring $\mathcal{O}(2^n)$ massive LLM forward passes, it is highly vulnerable to positional biases (e.g., the ``lost in the middle'' phenomenon), and it evaluates sequential orderings rather than permutation-invariant sets, fundamentally conflicting with the unordered nature of Shapley axioms.

To overcome these algorithmic bottlenecks, we introduce \textbf{Easton+}, an embedding-based formulation that completely decouples the heavy neural network inference from the exponential combinatorial evaluation. By mapping both the retrieved snippets and the generated response into a continuous, order-invariant dense vector space, Easton+ evaluates informational contribution through the rigorous lens of spatial geometry and linear algebra.

The primary contributions of this work are summarized as follows:
\begin{itemize}
    \item \textbf{Geometric Characteristic Functions:} We redefine the characteristic value function using the spatial volume of the parallelotope spanned by snippet embeddings, computed via the determinant of their Gram matrix. We introduce two variants---an additive difference ($v_{emb}$) and a multiplicative log-volume difference ($v_{\log}$)---which mathematically capture semantic novelty as the orthogonal distance between the model's response and the context subspace.
    \item \textbf{Theoretical Equivalence to LLM Likelihoods:} We formally bridge the geometric and probabilistic domains. We rigorously prove (Theorem \ref{thm:equi}) that if the LLM's conditional probability distribution satisfies the Determinantal Point Process (DPP) assumption, the Shapley values derived from our geometric log-volume function are strictly identical to those derived from the intractable LLM-based Easton algorithm.
    \item \textbf{Tractable, Closed-Form Marginal Gains:} We eliminate the need to re-evaluate the full characteristic function for every subset. By leveraging the Sherman-Morrison rank-1 update and orthogonal projection mechanics, we derive exact, closed-form marginal contributions. This breakthrough reduces the combinatorial evaluation bottleneck from $\mathcal{O}(2^n)$ GPU-bound LLM inferences to $\mathcal{O}(2^n \cdot d^2)$ CPU-bound algebraic operations, collapsing exact evaluation times from days to milliseconds and unlocking real-time polynomial-time approximations for massive context windows.
    \item \textbf{Spectral Entropy Generalization:} We extend the geometric framework to evaluate representational diversity via von Neumann spectral entropy. By defining a value function based on the divergence between the actual joint entropy and an orthogonal independent baseline, we formulate a measure of spectral mutual information that naturally penalizes semantic dilution and guarantees flawless boundary conditions.
\end{itemize}

By shifting the attribution paradigm from sequence-dependent likelihoods to space-spanning embedding volumes, Easton+ provides a theoretically grounded, fully scalable, and exact foundation for evaluating LLM context utilization.

\subsection{Related Work and Positioning}

The challenge of attributing Large Language Model (LLM) generation to source context within Retrieval-Augmented Generation (RAG) systems has sparked a line of research in perturbation-based post-hoc explanations. 
The canonical baseline, which we refer to as the log-likelihood Easton algorithm, was formalized by Nematov et al.~\cite{nematov2025source}. This approach frames source attribution as an $n$-player cooperative game, where the utility of a document subset $S'$ is evaluated using the conditional output log-likelihood $\log P(r \mid q, S')$ under teacher forcing~\cite{nematov2025source}. 
While theoretically elegant, exact evaluation remains computationally intractable due to the $\mathcal{O}(2^n)$ LLM forward passes, forcing practitioners to rely on linear surrogate approximations such as KernelSHAP~\cite{lundberg2017unified} or ContextCite~\cite{cohen2024contextcite}. 
Furthermore, these likelihood-based methods are severely distorted by sequential position bias (e.g., the ``lost in the middle'' phenomenon) and token-length artifacts~\cite{nematov2025source}.

In parallel, another line of research has leveraged spectral algebra to evaluate the representational diversity of continuous embeddings. 
Notably, the Vendi Score~\cite{friedman2022vendi} defines semantic diversity as the von Neumann spectral entropy of a normalized Gram similarity matrix, a formulation recently expanded to dataset valuation and scaling dynamics~\cite{bilmes2026dataset}. 
Our proposed framework, \textbf{Easton+}, bridges these previously disjoint paradigms. 
By mapping both context snippets and responses into a continuous, order-invariant vector space (or directly utilizing the pre-softmax attention Gram matrices of the Transformer), we replace the expensive, sequence-dependent probability evaluations of Nematov et al.~\cite{nematov2025source} with exact, closed-form geometric volume updates. 
This formulation not only reduces the combinatorial bottleneck from days of GPU compute to milliseconds of CPU algebra, but it also mathematically guarantees position-invariance and provides a rigorous spectral entropy-based formulation for multi-dimensional context attribution.


\section{Easton algorithm}

Let $p$ be a prompt consisting of system instructions, user context, and a query $q$. Let $S = [s_1, \dots, s_n]$ be the list of $n$ retrieved snippets. Let $r$ be the response generated using $p$ and $S$. For any order-preserving sublist $S' \subseteq S$, define its value function as the change in log-likelihood:
$$v(S') = \log P(r \mid p, S') - \log P(r \mid p)$$
(Note: $v(\emptyset) = 0$.)\\
{\bf Global Attribution:} Calculate the overall snippet weight $w$, bounded at zero:
$$w = \max\left(0, 1 - \frac{\log P(r \mid p, S)}{\log P(r \mid p)}\right)$$
The global contribution shares are then split between the model and the retrieved snippets. \\
{\bf Model (LLM) and snippets share:} We assume that the model contribution is $\text{contrib}(\text{model}) = 1 - w$  
and the rest corresponds to the snippets contribution, i.e. $\text{contrib}(S) = w$ \\
{\bf Marginal Contributions}: 
\begin{itemize}
    \item Shapley Values: For queries where the snippets improved the response ($w > 0$), calculate the Shapley value for each individual snippet $s_i$ (where $1 \le i \le n$):
    $$\text{shap}(i) = \sum_{S' \subseteq S \setminus \{s_i\}} \frac{\vert{}S'\vert{}!(n - \vert{}S'\vert{} - 1)!}{n!} \Big( v(S' \cup \{s_i\}) - v(S') \Big)$$
    (Note: The sum of all individual Shapley values perfectly reconstructs the total value: $\sum_{i=1}^n \text{shap}(i) = v(S)$)\\ 
    \item Normalization: Filter out any negative Shapley values and normalize them so they sum to $1$:$$\text{normshap}(i) = \frac{\max(0, \text{shap}(i))}{\sum_{j=1}^n \max(0, \text{shap}(j))}$$
    \item Final Attribution Distribute the global snippets share $w$ according to each snippet's normalized contribution:$$\text{contrib}(i) = w \cdot \text{normshap}(i)$$
\end{itemize}    
This ensures all attributions are properly scaled, guaranteeing the final property:$$\text{contrib}(\text{model}) + \sum_{i=1}^n \text{contrib}(i) = 1$$

\section{Limitations of the Easton Algorithm}
\begin{itemize}
    \item {\bf Intractable Computational Complexity}: Exact Shapley value attribution necessitates evaluating the log-likelihood for the power set of retrieved snippets, yielding an exponential $O(2^n)$ inference cost. This renders the approach computationally prohibitive for standard context windows (e.g., $n=20$ requires $>10^6$ model evaluations).
    \item {\bf Positional Bias and Sequence Dependence}: LLMs exhibit pronounced primacy and recency biases (the "lost in the middle" phenomenon). Consequently, the conditional likelihood is highly sensitive to snippet permutation, such that 
    $$P(r \mid [s_1, s_2]) \neq P(r \mid [s_2, s_1])$$
    This sequence dependence exacerbates as context length increases.
    \item {\bf Ill-Defined Characteristic Function:}
    Shapley values mathematically require a characteristic function $v(S)$ defined over an unordered set of elements. To rigorously define $v(S)$ for sequential LLM inputs, it must be formulated as the expected likelihood across all permutations $\Pi (S)$:
    $$v(S) = \mathbb{E}_{\pi \in \Pi(S)} [ P(r \mid \pi(S)) ]$$
    Computing this expectation introduces an additional $O(\vert S \vert !)$ factorial complexity. Conversely, evaluating the function using a static, default retriever ordering invalidates the mathematical foundation; it inextricably confounds the intrinsic informational value of a snippet with its arbitrary position within the prompt.
\end{itemize}


\section{Embedding-Based Easton: additive difference}

To address the positional bias inherent in LLM likelihood computations, we define an embedding-based variant of the Easton algorithm. Let $f$ be an embedding function (e.g., a standard dense retrieval model) that maps text to a vector space. We define the embedding vectors for the $n$ snippets as $E_i = f(s_i)$ for $1 \le i \le n$, and for the generated response as $R = f(r)$. For any subset of snippets $S' \subseteq S$, let $E_{S'}$ denote the corresponding set of embedding vectors. We assume that all embeddings are of unit length, such that $\|E_i\|_2 = 1$ for all $i$, and $\|R\|_2 = 1$.

\subsection{Volume of Spans}

To formally measure the span of our retrieval embeddings, we rely on the foundational geometric relationship between the volume of a parallelotope and the determinant of a Gram matrix.

\begin{lemma}\label{lemma:vol_det}
\textbf{(Volume and the Gram Determinant).} 
Let $E_{S'} = \{E_1, \dots, E_k\}$ be a set of $k$ linearly independent vectors in $\mathbb{R}^d$. The $k$-dimensional volume of the parallelotope spanned by these vectors is exactly equal to the square root of the determinant of their $k \times k$ Gram matrix $G$, where each entry is the inner product $G_{ij} = E_i \cdot E_j$.
\end{lemma}


The informational coverage of a snippet set can be quantified by the volume of the parallelotope formed by its embedding vectors. This volume is computed via the square root of the determinant of the Gram matrix $G$, where $G_{ij} = E_i \cdot E_j$:
$$ \text{Vol}(E_{S'}) = \sqrt{\det(G)} $$
Geometrically, if the vectors in $E_{S'}$ are highly redundant (semantically similar), the parallelotope flattens, driving the volume toward $0$. Conversely, if the vectors provide independent, orthogonal information, the volume is maximized.



\subsection{Value Function Based on Volume of Embeddings}

Using the spatial volume, we propose an embedding-based characteristic function $v_{emb}(S')$ that measures the informational contribution of a snippet subset $S'$. We define this as the difference in spatial volume before and after incorporating the response embedding $R$:
\begin{align} \label{eq:emb_val}
v_{emb}(S') = \text{Vol}(E_{S'}) - \text{Vol}(E_{S'} \cup \{R\})
\end{align}
Using the geometric properties of determinants, the volume of the augmented set $E_{S'} \cup \{R\}$ factors into the original volume multiplied by the orthogonal distance of $R$ to the span of $E_{S'}$:
$$ \text{Vol}(E_{S'} \cup \{R\}) = \text{Vol}(E_{S'}) \times d_{\perp}(R, \text{span}(E_{S'})) $$
Substituting this factorization back into Equation \ref{eq:emb_val} yields:
$$ v_{emb}(S') = \text{Vol}(E_{S'}) \times \big(1 - d_{\perp}(R, \text{span}(E_{S'}))\big) $$
The term $(1 - d_{\perp})$ functions as a scalar relevance score. It approaches $1$ if the response $R$ is perfectly spanned (predicted) by the snippets in $S'$, and $0$ if $R$ contains entirely novel, orthogonal information. Crucially, this relevance is scaled by $\text{Vol}(E_{S'})$, ensuring the model assigns higher characteristic values to snippet combinations that are highly diverse.


\subsubsection{Explanation of orthogonal distance}
To explicitly compute the orthogonal distance $d_{\perp}(R, \text{span}(E_{S'}))$, we rely on the linear algebraic projection of the response embedding $R$ onto the subspace spanned by the context embeddings $E_{S'}$. Let $E$ be a matrix whose columns are the embedding vectors in $S'$. The subspace generated by these snippets is the column space of $E$. 

Assuming the snippet vectors are linearly independent, we can construct the orthogonal projection matrix $P_{S'}$ using the Gram matrix $G = E^\top E$ defined previously:
$$ P_{S'} = E (E^\top E)^{-1} E^\top = E G^{-1} E^\top $$
This matrix maps any vector onto the span of $E_{S'}$. The projected response vector—representing the semantic information in $R$ that is entirely explained by the context $S'$—is given by $\hat{R} = P_{S'} R$. 

The strictly novel information contributed by the model response is isolated in the residual vector, $R_{\perp}$, which is orthogonal to the context subspace:
$$ R_{\perp} = R - \hat{R} = (I - P_{S'}) R $$
The orthogonal distance is then simply the Euclidean norm (magnitude) of this residual vector:
$$ d_{\perp}(R, \text{span}(E_{S'})) = \| R_{\perp} \|_2 $$
By integrating this back into our geometric volume factorization, we can see exactly how the residual orthogonal distance scales the characteristic value of the subset.

\begin{example}
\textbf{(2D Projection and Orthogonal Distance with Unit Vectors).}
Consider a simplified two-dimensional embedding space where a single retrieved snippet ($n=1$) forms our context subset $S' = \{s_1\}$, and we are evaluating a model response $r$. Assuming standard L2-normalized embeddings, let their respective unit vectors ($\|E_1\| = 1, \|R\| = 1$) be:
$$ E_1 = \begin{bmatrix} 1 \\ 0 \end{bmatrix}, \quad R = \begin{bmatrix} 0.6 \\ 0.8 \end{bmatrix} $$
The subspace $\text{span}(E_1)$ is simply the one-dimensional line along the x-axis. We first compute the orthogonal projection of $R$ onto $E_1$. Since $E$ is a single column vector, the matrix projection formula simplifies to the standard vector scalar projection:
$$ \hat{R} = \frac{R \cdot E_1}{\|E_1\|^2} E_1 = \frac{(0.6)(1) + (0.8)(0)}{1^2 + 0^2} \begin{bmatrix} 1 \\ 0 \end{bmatrix} = \frac{0.6}{1} \begin{bmatrix} 1 \\ 0 \end{bmatrix} = \begin{bmatrix} 0.6 \\ 0 \end{bmatrix} $$
The residual vector, representing the orthogonal novelty of the model response, is the difference between the original vector and its projection:
$$ R_{\perp} = R - \hat{R} = \begin{bmatrix} 0.6 \\ 0.8 \end{bmatrix} - \begin{bmatrix} 0.6 \\ 0 \end{bmatrix} = \begin{bmatrix} 0 \\ 0.8 \end{bmatrix} $$
Finally, the orthogonal distance is the magnitude of this residual:
$$ d_{\perp}(R, \text{span}(E_1)) = \|R_{\perp}\|_2 = \sqrt{0^2 + 0.8^2} = 0.8 $$
\end{example}

\subsection{Exact Marginal Contributions and Guaranteed Non-Negativity}

To compute the exact Shapley values for the embedding-based Easton algorithm, we must evaluate the marginal contribution of adding a specific snippet $s_i$ to an existing context subset $S'$. Using the regularized information gain characteristic function, this marginal value is defined as:
\begin{align}\label{eq:v_emp}
 \Delta v_{emb}(S', s_i) = v_{emb}(S' \cup \{s_i\}) - v_{emb}(S') \enspace .
\end{align}
Instead of performing computationally expensive matrix inversions for every combinatorial subset, we compute this marginal gain directly in closed form using the Sherman-Morrison rank-1 update formula. Let $M = A_{S'}^{-1}$ be the inverse regularized covariance matrix of the current subset. The exact algebraic marginal gain simplifies to:
\begin{align}  \label{eq:marginal_gain}
  \Delta v_{emb}(S', s_i) = - \log \left( 1 - \frac{(R^\top M E_i)^2}{(1 + R^\top M R)(1 + E_i^\top M E_i)} \right)
\end{align}
Crucially, this formulation mathematically guarantees that the marginal contribution is \textbf{strictly non-negative} ($\Delta v_{emb}(S', s_i) \ge 0$). Because the regularized inverse matrix $M$ is symmetric and positive definite, it defines a valid inner product space. Consequently, the Cauchy-Schwarz inequality mandates that the squared cross-term is bounded by the product of the individual norms:
$$ (R^\top M E_i)^2 \le (R^\top M R)(E_i^\top M E_i) $$
This bounds the fractional term inside the equation entirely within the interval $[0, 1)$. Subtracting this fraction from $1$ guarantees a scalar strictly greater than $0$ and less than or equal to $1$. Because the natural logarithm of any value in $(0, 1]$ is negative or zero, the leading negative sign permanently inverts the result, ensuring the final information gain is always positive or zero. This geometric property mirrors the principle of monotonicity: adding new spatial dimensions to the context subspace can never decrease its overall capacity to span the model's response.

\subsection{Information Theoretic interpretation of $\Delta v_{emb}(S', s_i)$ (optional, to be revised)}

From an information-theoretic perspective, the marginal value $\Delta v_{emb}(S', s_i)$ formally quantifies the conditional information gain—mathematically analogous to conditional mutual information within a Gaussian channel—that snippet $s_i$ provides regarding the response $R$, strictly conditioned on the prior context $S'$. To elucidate this connection, we can model the continuous embedding space as a statistical manifold where the context subset parameterizes a zero-mean multivariate Gaussian distribution, $\mathcal{N}(0, A_{S'})$. Under this distribution, the ``surprisal'' or novelty of the response embedding $R$ given the context is explicitly defined by the monotonic function:
$$ N(S') = \log(1 + R^\top A_{S'}^{-1} R) $$
The core of this novelty metric is the quadratic form $R^\top A_{S'}^{-1} R$, which is the squared Mahalanobis distance. In our Gaussian statistical model, this distance is directly proportional to the negative log-likelihood of observing the target response $R$ given the context covariance $A_{S'}$. Thus, $N(S')$ rigorously measures how ``surprising'' or unpredicted the model response is relative to the information already present in the snippet subset.

In the information geometry of this space, the precision matrix $M = A_{S'}^{-1}$ is precisely the Fisher Information Matrix (FIM), acting as the local Riemannian metric tensor that defines the curvature of the semantic manifold. Consequently, the critical cross-term $R^\top M E_i$ governing the marginal gain evaluates the exact Fisher inner product, $\langle R, E_i \rangle_M$. Within this Riemannian geometry, the vector $M E_i$ corresponds to the natural gradient of the Kullback-Leibler (KL) divergence. The Fisher inner product therefore measures the exact magnitude by which the new snippet $E_i$ updates the context's predictive log-likelihood toward the target response $R$, dynamically scaled by the statistical precision of the semantic information already spanned by $S'$.

Furthermore, this geometric structure establishes a profound mathematical equivalence between LLM attribution and sequential decision-making, specifically the exploration mechanisms in linear contextual bandits (e.g., LinUCB) and active learning. In the contextual bandit framework, the regularized covariance matrix $A_{S'}$ acts as the design matrix, summarizing the features of previously pulled arms (the context). To minimize predictive variance, these algorithms select the most informative arms by maximizing an upper confidence bound strictly proportional to the Mahalanobis distance, $E_i^\top A_{S'}^{-1} E_i$. By evaluating the marginal gain $\Delta v_{emb}$ via these precise precision-matrix updates, our formulation inherently mirrors this principle of information maximization. The characteristic function scores the marginal attribution of snippet $s_i$ exactly as a bandit algorithm would value an arm that maximally resolves the remaining semantic uncertainty regarding the target reward state, $R$.


%%%%%%%%%%%%%%%%%%%%%%%%%%%%%%%%%%%%%%%%%%%%%%%%%%%%%%%%%%%%%%%%%%%%%%%%%%%%%%%%%%%%%%%%%%%%%%%%%%%%%%%%%%%%
%%%%%%%%%%%%%%%%%%%%%%%%%%%%%%%%%%%%%%%%%%%%%%%%%%%%%%%%%%%%%%%%%%%%%%%%%%%%%%%%%%%%%%%%%%%%%%%%%%%%%%%%%%%%

\section{Embedding-Based Easton: multiplicative difference}

For the Shapley values derived from $v_{emb}(S')$ to exactly match those derived from the original log-likelihood formulation $v(S') = \log P(r \mid p, S') - \log P(r \mid p)$, the following conditions must be satisfied:

% \begin{condition}
% \textbf{The Determinantal Point Process (DPP) Assumption.} 
% The language model's underlying probability distribution must be representable as a DPP. In a DPP, the probability of observing a subset of items is directly proportional to the squared volume of the parallelepiped they span:
% $$ P(r \mid S') \propto \text{Vol}(E_{S'} \cup \{R\})^2 $$
% \end{condition}

% \begin{remark}
% The embedding models are typically trained by using triplet loss, at least Gecko is trained like that. Whereas $v(S')$ is some kind of conditional probability. We need to test to what extent this Condition A is true for various queries.
% \end{remark}

% \begin{condition}
% \textbf{A Logarithmic Link Function.} 
% The LLM value function is logarithmic, whereas the proposed volume difference is linear. Because $\log P(r \mid p, S') \propto \log P(r, S') - \log P(S')$, the geometric approximation must be modified from a raw difference of volumes to a difference of log-volumes:
% $$ v_{\log}(S') = \log(\text{Vol}(E_{S'} \cup \{R\})) - \log(\text{Vol}(E_{S'})) $$
% Given the orthogonal distance property defined earlier, this simplifies to the log of the orthogonal distance:
% $$ v_{\log}(S') = \log\big(d_{\perp}(R, \text{span}(E_{S'}))\big) $$
% \end{condition}









% To demonstrate that the embedding-based value function yields identical attributions to the original log-likelihood Easton algorithm, we must prove that their normalized Shapley values are equivalent under the Determinantal Point Process (DPP) assumption (Condition A) and the logarithmic link (Condition B).

% \begin{theorem}\label{thm:equi}
% \textbf{(Equivalence of LLM and Geometric Shapley Values).} 
% \textit{Let $S = \{s_1, \dots, s_n\}$ be a set of retrieved snippets and $r$ be a generated response, with corresponding embedding vectors $E_S$ and $R$. Define the log-likelihood characteristic function as $v_{LL}(S') = \log P(r \mid S') - \log P(r \mid \emptyset)$ and the geometric characteristic function as $v_{\log}(S') = \log \text{Vol}(E_{S'} \cup \{R\}) - \log \text{Vol}(E_{S'})$. If the language model's joint probability distribution satisfies the Determinantal Point Process (DPP) assumption such that $P(r, S') = K \cdot \text{Vol}(E_{S'} \cup \{R\})^2$ and $P(S') = K \cdot \text{Vol}(E_{S'})^2$ for some constant $K$, then the normalized Shapley values derived from $v_{LL}$ and $v_{\log}$ are strictly identical for all $s_i \in S$.}
% \end{theorem}
% \begin{proofof}{Theorem \ref{thm:equi}}
% By definition, the conditional probability of the response $r$ given a snippet subset $S'$ is $P(r \mid S') = \frac{P(r, S')}{P(S')}$. Applying the given DPP assumption, this ratio becomes:
% $$ P(r \mid S') = \frac{K \cdot \text{Vol}(E_{S'} \cup \{R\})^2}{K \cdot \text{Vol}(E_{S'})^2} = \left( \frac{\text{Vol}(E_{S'} \cup \{R\})}{\text{Vol}(E_{S'})} \right)^2 $$
% Taking the logarithm of both sides transforms the squared ratio into a linear difference of log-volumes:
% \begin{align*}
% \log P(r \mid S') &= 2 \Big( \log \text{Vol}(E_{S'} \cup \{R\}) - \log \text{Vol}(E_{S'}) \Big) \\
% &= 2 \cdot v_{\log}(S')
% \end{align*}
% The marginal contribution of adding a snippet $s_i$ to subset $S'$ in the original Easton characteristic function is defined as $\Delta v_{LL}(S', s_i) = v_{LL}(S' \cup \{s_i\}) - v_{LL}(S')$. Substituting our geometric equivalence into this definition reveals that the baseline constant $2 \cdot v_{\log}(\emptyset)$ cancels out:
% \begin{align*}
% \Delta v_{LL}(S', s_i) &= \Big[ 2 v_{\log}(S' \cup \{s_i\}) - 2 v_{\log}(\emptyset) \Big] - \Big[ 2 v_{\log}(S') - 2 v_{\log}(\emptyset) \Big] \\
% &= 2 \Big( v_{\log}(S' \cup \{s_i\}) - v_{\log}(S') \Big) \\
% &= 2 \cdot \Delta v_{\log}(S', s_i)
% \end{align*}
% Because the Shapley value operator $\Phi$ is mathematically linear, scaling all marginal contributions by a constant factor of $2$ simply scales the resulting Shapley value by the same factor:
% $$ \text{shap}_{LL}(i) = 2 \cdot \text{shap}_{\log}(i) $$
% The final step of the Easton algorithm normalizes these values. Let $\text{normshap}_{LL}(i)$ be the normalized Shapley value for snippet $s_i$:
% $$ \text{normshap}_{LL}(i) = \frac{\max(0, 2 \cdot \text{shap}_{\log}(i))}{\sum_{j=1}^n \max(0, 2 \cdot \text{shap}_{\log}(j))} $$
% Because the scalar factor $2$ is strictly positive, it commutes with the $\max(0, \cdot)$ operator and factors out of both the numerator and the summation in the denominator. The scalar cancels completely, yielding:
% $$ \text{normshap}_{LL}(i) = \text{normshap}_{\log}(i) $$
% This completes the proof.
% \end{proofof}


\begin{condition}
\textbf{The Determinantal Point Process (DPP) Assumption.} 
The language model's underlying probability distribution, conditioned on the prompt $p$, must be representable as a DPP. In a DPP, the probability of observing a subset of items is directly proportional to the squared volume of the parallelepiped they span:
$$ P(r, S' \mid p) \propto \text{Vol}(E_{S'} \cup \{R\})^2 $$
\end{condition}

\begin{remark}
The embedding models are typically trained by using triplet loss, at least Gecko is trained like that. Whereas $v(S')$ is some kind of conditional probability. We need to test to what extent this Condition A is true for various queries.
\end{remark}

\begin{condition}
\textbf{A Logarithmic Link Function.} 
The LLM value function is logarithmic, whereas the proposed volume difference is linear. Because $\log P(r \mid p, S') \propto \log P(r, S' \mid p) - \log P(S' \mid p)$, the geometric approximation must be modified from a raw difference of volumes to a difference of log-volumes:
$$ v_{\log}(S') = \log(\text{Vol}(E_{S'} \cup \{R\})) - \log(\text{Vol}(E_{S'})) $$
Given the orthogonal distance property defined earlier, this simplifies to the log of the orthogonal distance:
\begin{align} \label{eq:v_log}
     v_{\log}(S') = \log\big(d_{\perp}(R, \text{span}(E_{S'}))\big)
\end{align}     
\end{condition}

To demonstrate that the embedding-based value function yields identical attributions to the original log-likelihood Easton algorithm, we must prove that their normalized Shapley values are equivalent under the Determinantal Point Process (DPP) assumption (Condition A) and the logarithmic link (Condition B).

\begin{theorem}\label{thm:equi}
\textbf{(Equivalence of LLM and Geometric Shapley Values).} 
\textit{Let $p$ be a prompt, $S = \{s_1, \dots, s_n\}$ be a set of retrieved snippets, and $r$ be a generated response, with corresponding embedding vectors $E_S$ and $R$. Define the log-likelihood characteristic function as $v_{LL}(S') = \log P(r \mid p, S') - \log P(r \mid p, \emptyset)$ and the geometric characteristic function as $v_{\log}(S') = \log \text{Vol}(E_{S'} \cup \{R\}) - \log \text{Vol}(E_{S'})$. If the language model's conditional probability distribution satisfies the Determinantal Point Process (DPP) assumption such that $P(r, S' \mid p) = K_p \cdot \text{Vol}(E_{S'} \cup \{R\})^2$ and $P(S' \mid p) = K_p \cdot \text{Vol}(E_{S'})^2$ for some prompt-dependent constant $K_p$, then the normalized Shapley values derived from $v_{LL}$ and $v_{\log}$ are strictly identical for all $s_i \in S$.}
\end{theorem}
\begin{proofof}{Theorem \ref{thm:equi}}
By definition, the conditional probability of the response $r$ given the prompt $p$ and snippet subset $S'$ is $P(r \mid p, S') = \frac{P(r, S' \mid p)}{P(S' \mid p)}$. Applying the given DPP assumption, this ratio becomes:
$$ P(r \mid p, S') = \frac{K_p \cdot \text{Vol}(E_{S'} \cup \{R\})^2}{K_p \cdot \text{Vol}(E_{S'})^2} = \left( \frac{\text{Vol}(E_{S'} \cup \{R\})}{\text{Vol}(E_{S'})} \right)^2 $$
Taking the logarithm of both sides transforms the squared ratio into a linear difference of log-volumes:
\begin{align*}
\log P(r \mid p, S') &= 2 \Big( \log \text{Vol}(E_{S'} \cup \{R\}) - \log \text{Vol}(E_{S'}) \Big) \\
&= 2 \cdot v_{\log}(S')
\end{align*}
The marginal contribution of adding a snippet $s_i$ to subset $S'$ in the original Easton characteristic function is defined as $\Delta v_{LL}(S', s_i) = v_{LL}(S' \cup \{s_i\}) - v_{LL}(S')$. Substituting our geometric equivalence into this definition reveals that the baseline constant $2 \cdot v_{\log}(\emptyset)$ cancels out:
\begin{align*}
\Delta v_{LL}(S', s_i) &= \Big[ 2 v_{\log}(S' \cup \{s_i\}) - 2 v_{\log}(\emptyset) \Big] - \Big[ 2 v_{\log}(S') - 2 v_{\log}(\emptyset) \Big] \\
&= 2 \Big( v_{\log}(S' \cup \{s_i\}) - v_{\log}(S') \Big) \\
&= 2 \cdot \Delta v_{\log}(S', s_i)
\end{align*}
Because the Shapley value operator $\Phi$ is mathematically linear, scaling all marginal contributions by a constant factor of $2$ simply scales the resulting Shapley value by the same factor:
$$ \text{shap}_{LL}(i) = 2 \cdot \text{shap}_{\log}(i) $$
The final step of the Easton algorithm normalizes these values. Let $\text{normshap}_{LL}(i)$ be the normalized Shapley value for snippet $s_i$:
$$ \text{normshap}_{LL}(i) = \frac{\max(0, 2 \cdot \text{shap}_{\log}(i))}{\sum_{j=1}^n \max(0, 2 \cdot \text{shap}_{\log}(j))} $$
Because the scalar factor $2$ is strictly positive, it commutes with the $\max(0, \cdot)$ operator and factors out of both the numerator and the summation in the denominator. The scalar cancels completely, yielding:
$$ \text{normshap}_{LL}(i) = \text{normshap}_{\log}(i) $$
This completes the proof.
\end{proofof}

\subsection{Exact Marginal Contributions for the Geometric Log-Volume}

To compute the exact Shapley values for the pure geometric characteristic function $v_{\log}$, we must evaluate the marginal change in the log-volume when a specific snippet $s_i$ is added to an existing context subset $S'$. Based on the logarithmic link (Condition B), this marginal value evaluates the change in orthogonal distance:
$$ \Delta v_{\log}(S', s_i) = v_{\log}(S' \cup \{s_i\}) - v_{\log}(S') $$
which is
$$ \Delta v_{\log}(S', s_i) = \log\big(d_{\perp}(R, \text{span}(E_{S'} \cup \{E_i\}))\big) - \log\big(d_{\perp}(R, \text{span}(E_{S'}))\big) $$
Instead of computing the determinants of entirely new Gram matrices for every subset, we can evaluate this update geometrically. Let $R_{\perp}$ denote the residual of the response vector after projecting out the span of $S'$, such that $\|R_{\perp}\| = d_{\perp}(R, \text{span}(E_{S'}))$. Similarly, let $E_{i, \perp}$ be the residual of the new snippet embedding $E_i$ orthogonal to $S'$. 

Using standard projection mechanics, the new squared orthogonal distance after adding $E_i$ updates as:
$$ d_{\perp}^2(R, \text{span}(E_{S'} \cup \{E_i\})) = \|R_{\perp}\|^2 - \frac{(R_{\perp} \cdot E_{i, \perp})^2}{\|E_{i, \perp}\|^2} $$

By factoring out the original squared distance $\|R_{\perp}\|^2$ and applying the logarithmic property $\log(x) = \frac{1}{2} \log(x^2)$, the exact closed-form marginal gain elegantly simplifies to:
\begin{align} \label{eq:marginal_vlog}
\Delta v_{\log}(S', s_i) &= \frac{1}{2} \log \left( \frac{\|R_{\perp}\|^2 - \frac{(R_{\perp} \cdot E_{i, \perp})^2}{\|E_{i, \perp}\|^2}}{\|R_{\perp}\|^2} \right) \nonumber \\
&= \frac{1}{2} \log \left( 1 - \frac{(R_{\perp} \cdot E_{i, \perp})^2}{\|R_{\perp}\|^2 \|E_{i, \perp}\|^2} \right)
\end{align}

Crucially, this formulation reveals the exact mathematical bounds of the geometric approach. The fractional term inside the logarithm is the squared cosine of the angle between the residual vectors, $\cos^2(\theta)$. By the Cauchy-Schwarz inequality, this fraction is strictly bounded within $[0, 1]$. 

Because subtracting any value in $[0, 1]$ from $1$ yields a scalar $\le 1$, and the natural logarithm of any number $\le 1$ is negative or zero, we prove that \textbf{$\Delta v_{\log}(S', s_i)$ is strictly non-positive} ($\Delta v_{\log} \le 0$). 

This mathematically encodes the principle of geometric monotonicity: adding a new snippet dimension to the context subspace can only \textit{reduce} (or leave unchanged) the orthogonal distance of the response. Unlike $v_{emb}$, which represents an information \textit{gain} (and is thus positive), $v_{\log}$ fundamentally tracks distance \textit{reduction}.


\section{The Explicit Embedding-Based Easton Algorithm}

With the marginal gains derived in closed form, the exact Shapley computation over the $O(2^n)$ subsets becomes highly optimized. The full explicit algorithm operates as follows where $\mathcal{M} \in \{\text{emp}, \text{log} \}$:

\begin{itemize}
    \item \textbf{Step 1: Global Attribution (LLM Space).} \\
    We preserve the exact true-likelihood scaling of the original algorithm by anchoring the total snippet weight $w$ with exactly \textit{two} LLM forward passes (using the full context $S$ and empty context $\emptyset$):
    $$w = \max\left(0, 1 - \frac{\log P(r \mid p, S)}{\log P(r \mid p)}\right)$$
    The model share is fixed as $\text{contrib}(\text{model}) = 1 - w$.
    
    \item \textbf{Step 2: Exact Marginal Contributions (Embedding Space).} \\
    For queries where $w > 0$, compute the Shapley values geometrically. Using Equation \eqref{eq:marginal_gain} or \eqref{eq:marginal_vlog}, compute:
    $$ \text{shap}_{\mathcal{M}}(i) = \sum_{S' \subseteq S \setminus \{s_i\}} \frac{\vert S' \vert ! (n - \vert S' \vert - 1)!}{n!} \Delta v_{\mathcal{M}}(S', s_i) $$
    \textit{Computational Note:} By traversing the subsets using a Gray code order, the inverse matrix $M$ can be maintained sequentially via $O(d^2)$ updates. The entire summation evaluates in milliseconds on standard CPU hardware, entirely bypassing the exponential LLM bottleneck.
    
    \item \textbf{Step 3: Normalization.} \\
    Because Equation \ref{eq:marginal_gain} guarantees non-negativity, the $\max(0, \cdot)$ filter from the original Easton algorithm is structurally redundant. We safely normalize the values so they sum to $1$:
    $$ \text{normshap}_{\mathcal{M}}(i) = \frac{\text{shap}_{\mathcal{M}}(i)}{\sum_{j=1}^n \text{shap}_{\mathcal{M}}(j)} $$
    
    \item \textbf{Step 4: Final Attribution.} \\
    Distribute the global snippets share $w$ according to each snippet's normalized geometric contribution:
    $$ \text{contrib}(i) = w \cdot \text{normshap}_{\mathcal{M}}(i) $$
\end{itemize}


\subsection{Computational Complexity Analysis}

To be mathematically precise, the exact computation of Shapley values inherently requires evaluating marginal contributions over the power set of $n$ items. Therefore, the combinatorial complexity remains $\mathcal{O}(2^n)$ in both formulations. The fundamental breakthrough of the embedding-based approach lies in decoupling the neural network inference from this exponential combinatorial loop, shifting the asymptotic computational bottleneck.

\textbf{LLM-Based Easton Complexity (Intractable).} 
In the original Easton algorithm, evaluating a single subset requires a full autoregressive LLM forward pass to calculate the conditional likelihood $P(r \mid p, S')$. Let $\mathcal{C}_{LLM}$ denote the massive floating-point cost of this transformer inference. The total complexity for exact attribution is:
$$ \mathcal{O}(2^n \cdot \mathcal{C}_{LLM}) $$
For a standard retrieval context of $n = 20$ snippets, evaluating the $1.04 \times 10^6$ combinations requires an intractable amount of GPU time per query, rendering it dead on arrival for production systems.

\textbf{Embedding-Based Easton Complexity (Exact Computation).} 
The geometric formulation fully decouples the heavy neural inference from the combinatorial sum, splitting the computation into two phases:
\begin{itemize}
    \item \textbf{Embedding Phase:} Dense vectors for the $n$ snippets and the response $r$ are computed exactly once. The complexity is $\mathcal{O}(n \cdot \mathcal{C}_{emb})$, where $\mathcal{C}_{emb}$ is a single pass through a lightweight encoder model.
    \item \textbf{Attribution Phase:} Using a Gray code sequence to traverse the power set, each successive subset differs by exactly one snippet (either an addition or removal). By applying the Sherman-Morrison rank-1 update, the inverse covariance matrix and the subsequent marginal contribution are calculated in $\mathcal{O}(d^2)$ algebraic operations, where $d$ is the embedding dimension. 
\end{itemize}
The total exact complexity becomes:
$$ \mathcal{O}(n \cdot \mathcal{C}_{emb} + 2^n \cdot d^2) $$
Because $\mathcal{O}(d^2)$ vector operations execute in microseconds on a standard CPU, exact evaluation for standard context lengths ($n \le 20$) drops from days of GPU compute to milliseconds of CPU compute.

\textbf{Polynomial-Time Approximation for Large Contexts ($n \gg 20$).} 
For extended context windows (e.g., $n = 50$ or $100$), computing $2^n$ evaluations becomes physically impossible regardless of the per-step efficiency. However, the ultra-fast $\mathcal{O}(d^2)$ algebraic evaluation natively unlocks Monte Carlo Shapley Approximation. 

By sampling $K$ random permutations of the snippets and computing the marginal gains along each path, the true Shapley values can be approximated with strict theoretical guarantees. The combinatorial complexity drops to strictly polynomial time:
$$ \mathcal{O}(n \cdot \mathcal{C}_{emb} + K \cdot n \cdot d^2) $$
While evaluating thousands of permutation paths remains prohibitively slow for LLMs ($\mathcal{O}(K \cdot n \cdot \mathcal{C}_{LLM})$), it executes near-instantaneously in the embedding space. This guarantees scalable, real-time attribution for virtually infinitely large context windows.

\section{Spectral Interpretation and an Alternative Value Function}

The log-volume value function defined in Equation \eqref{eq:v_log} possesses a natural spectral interpretation. Recalling the fundamental matrix identity $$\log\det(A) = \text{Tr}(\log A)\enspace ,$$ 
the spatial log-volume is mathematically equivalent to the trace of the matrix logarithm. In terms of the matrix's eigenvalues $\lambda_i$, this corresponds to the unnormalized sum of their logarithms: $\sum_i \log \lambda_i$. Consequently, the characteristic function $v_{emb}(S')$ evaluates information gain by taking the difference of the trace of the log-covariance matrices before and after incorporating the response embedding $R$.

For the geometric value function $v_{\log}(S')$ defined in Condition B, let $G_{S'}$ denote the Gram matrix of the context embeddings. Using the trace operator and the eigenvalues $\lambda_i$, the log-volume difference can be explicitly written as:
\begin{align} \label{eq:v_log_trace}
v_{\log}(S') &= \frac{1}{2} \Big( \text{Tr}\big(\log G_{S' \cup \{R\}}\big) - \text{Tr}\big(\log G_{S'}\big) \Big) \nonumber \\
&= \frac{1}{2} \left( \sum_{i=1}^{\vert S' \vert + 1} \log \lambda_i(G_{S' \cup \{R\}}) - \sum_{i=1}^{\vert S' \vert} \log \lambda_i(G_{S'}) \right)
\end{align}

Similarly, the regularized information gain $v_{emb}(S')$ defined in Equation \eqref{eq:emb_val} operates on the $d \times d$ covariance matrices $A_{S'}$. TODO: work out the spectral interpretation

% The conditional surprisal $N(S')$ introduced by the response $R$ is explicitly the change in the trace of the log-covariance:
% $$ N(S') = \text{Tr}\big(\log(A_{S'} + R R^\top)\big) - \text{Tr}\big(\log A_{S'}\big) = \sum_{j=1}^d \log \lambda_j(A_{S'} + R R^\top) - \sum_{j=1}^d \log \lambda_j(A_{S'}) $$
% Consequently, the characteristic function $v_{emb}(S') = N(\emptyset) - N(S')$ evaluates to the difference of these unnormalized sum-of-log spectra:
% \begin{align} \label{eq:v_emb_trace}
% v_{emb}(S') &= \Big( \text{Tr}\big(\log(A_{\emptyset} + R R^\top)\big) - \text{Tr}\big(\log A_{\emptyset}\big) \Big) - \Big( \text{Tr}\big(\log(A_{S'} + R R^\top)\big) - \text{Tr}\big(\log A_{S'}\big) \Big) \nonumber \\
% &= \sum_{j=1}^d \log \left( \frac{\lambda_j(A_{\emptyset} + R R^\top)}{\lambda_j(A_{\emptyset})} \right) - \sum_{j=1}^d \log \left( \frac{\lambda_j(A_{S'} + R R^\top)}{\lambda_j(A_{S'})} \right)
% \end{align}
% In both formulations, information gain inherently evaluates spatial diversity via the \textit{unnormalized log-spectrum} of the system's similarity or covariance matrices.



An alternative, theoretically principled approach to quantifying representational diversity is the spectral entropy \cite{friedman2022vendi}. Rather than computing the trace of the log-spectrum, the spectral entropy evaluates the \textit{entropy} of the normalized eigenvalues. For a subset of embeddings with a Gram similarity matrix $K_{S'}$, one first defines a density matrix 
$$\rho_{S'} = \frac{G_{S'}}{\text{Tr}(G_{S'})}$$ 
such that its eigenvalues $p_i$ form a valid probability distribution ($\sum_i p_i = 1$). The diversity of the subset is then measured using the von Neumann entropy, 
$$\mathcal{H}(\rho_{S'}) = - \sum_i p_i \log p_i \enspace .$$ 

Using this spectral entropy framework, we can formulate an alternative  function, $v_{ent}(S')$. Analogous to the original geometric surprisal, we define the conditional novelty introduced by the response $R$ given context $S'$ as the marginal increase in the spectral entropy:
$$ v_{ent}(S') = \mathcal{H}(\rho_{S' \cup \{R\}}) - \mathcal{H}(\rho_{S'}) $$


Because spectral entropy strictly decreases (or resists growth) when the response $R$ is semantically spanned by the context $S'$, we can frame our characteristic value function around this exact behavior. We require the value function to be maximized when the entropy drops (indicating the context perfectly explains the response) and to evaluate to zero when $R$ is entirely novel.

However, because the density matrix $\rho$ normalizes by the trace (which inherently grows from $\vert S' \vert$ to $\vert S' \vert + 1$ when appending $R$), directly subtracting the prior entropy from the joint entropy introduces a scaling artifact. To eliminate this and maintain strict Shapley boundary conditions, we define the value function as the reduction in entropy relative to an assumption of perfect independence.

Let $\mathcal{H}_{ortho}(S' \cup \{R\})$ be the expected spectral entropy calculated under the assumption that $R$ is perfectly orthogonal to the entire subspace of $S'$. In this independent state, the response adds a completely novel, non-overlapping dimension to the system. We define the entropy-based characteristic function, $v_{ent}(S')$, as the divergence between this orthogonal expectation and the actual joint entropy:
\begin{align} \label{eq:v_ent_mi}
v_{ent}(S') = \mathcal{H}_{ortho}(S' \cup \{R\}) - \mathcal{H}(S' \cup \{R\})
\end{align}

This formulation acts as a form of spectral mutual information and guarantees flawless geometric boundary conditions:
\begin{itemize}
    \item \textbf{Complete Orthogonality ($R$ is novel):} If the response shares no semantic overlap with the retrieved context, the actual joint entropy perfectly matches the orthogonal baseline. The subtraction perfectly cancels out, yielding $v_{ent}(S') = 0$. The snippet subset rightfully receives no characteristic value.
    \item \textbf{Complete Span ($R$ is explained):} If the response is fully spanned by the retrieved context, the actual joint entropy shrinks as the probability mass concentrates on the pre-existing semantic dimensions. Because the actual entropy drops significantly below the orthogonal expectation, the difference is strictly positive, maximizing $v_{ent}(S')$.
\end{itemize}
By measuring how much the context subset restricts the semantic entropy of the response from reaching its orthogonal maximum, $v_{ent}(S')$ rigorously quantifies true information gain.

\subsection{TODO: entropy submodularity}


\section{Stochastic Batch Updates via Marginal Gain Sampling (WIP, does not make sense in this form)}

In practice, computing the exact raw Shapley value $\phi_i$ to serve as a static ground-truth label $y_i$ requires $\mathcal{O}(2^n)$ evaluations per snippet before training even begins. However, we can bypass this static pre-computation entirely by intertwining the model's gradient descent with Monte Carlo Shapley sampling.

Mathematically, the raw Shapley value is the expected value of the marginal gains over a specific subset distribution. Let $\mathcal{P}(S')$ define the discrete Shapley probability measure over all subsets $S' \subseteq S \setminus \{s_i\}$, where the probability of sampling a specific subset is inversely proportional to the number of subsets of that size:
$$ \mathcal{P}(S') = \frac{1}{n} \binom{n-1}{\vert S' \vert}^{-1} $$
The exact Shapley value is simply this expectation:
$$ \phi_i = \mathbb{E}_{S' \sim \mathcal{P}} \big[ \Delta v(S', s_i) \big] $$

A fundamental property of the squared error loss is that the minimizer of $\mathbb{E}[(X - c)^2]$ is exactly the mean, $c = \mathbb{E}[X]$. Therefore, instead of training the model $f_\theta(E_i, R)$ to regress on a pre-computed $\phi_i$, we define a \textbf{stochastic marginal loss} that evaluates the prediction directly against a sampled marginal gain:
$$ \mathcal{L}_{stochastic}(\theta) = \mathbb{E}_{S' \sim \mathcal{P}} \left[ \frac{1}{2} \Big( \Delta v(S', s_i) - f_\theta(E_i, R) \Big)^2 \right] $$
By taking the derivative with respect to the network parameters $\theta$, we obtain a gradient that is explicitly driven by the marginal gain of the subset:
$$ \nabla_\theta \mathcal{L}_{stochastic} = - \mathbb{E}_{S' \sim \mathcal{P}} \left[ \Big( \Delta v(S', s_i) - f_\theta(E_i, R) \Big) \nabla_\theta f_\theta(E_i, R) \right] $$

\paragraph{The Stochastic Batch Update Rule.}
This formulation yields a highly efficient batch update mechanism for training. During each training iteration, rather than requiring the exact $\phi_i$, the system executes the following steps for a batch of queries:
\begin{enumerate}
    \item Sample a context subset $S' \sim \mathcal{P}$ for a given target snippet $s_i$.
    \item Compute the exact closed-form marginal gain $\Delta v_{emb}(S', s_i)$ using the Sherman-Morrison update (Equation \ref{eq:marginal_gain}).
    \item Forward-pass the pointwise network to obtain the current prediction $f_\theta(E_i, R)$.
    \item Compute the stochastic residual error: $\epsilon = \Delta v_{emb}(S', s_i) - f_\theta(E_i, R)$.
    \item Augment the model weights using the backpropagated gradient: 
    $$ \theta \leftarrow \theta + \alpha \cdot \epsilon \cdot \nabla_\theta f_\theta(E_i, R) $$
\end{enumerate}

This mechanism transforms the exponential combinatorial burden of Shapley values into a standard stochastic learning process. The sampled marginal gains dynamically augment the gradient magnitude and direction. As training progresses over sufficiently large datasets and epochs, the network $f_\theta(E_i, R)$ asymptotically converges to predicting the true expected marginal contribution—the exact raw Shapley value—without the system ever computing an $\mathcal{O}(2^n)$ summation.


\section{Listwise Prediction and Consistent Empirical Risk Minimization}

The pointwise formulation intrinsically assumes that the attribution of snippet $s_i$ can be predicted independently of the other snippets in the context. However, the fundamental nature of the Shapley value is combinatorial and context-dependent; by the efficiency axiom, the sum of all individual attributions must perfectly equal the total value of the set, $\sum_{i=1}^n \phi_i = v(S)$. To rigorously model this joint distribution, we must transition from a pointwise architecture to a \textit{listwise} learning-to-rank paradigm, where the model evaluates the entire set of snippets simultaneously.

Let $E_S = [E_1, \dots, E_n]^\top \in \mathbb{R}^{n \times d}$ be the matrix of all snippet embeddings for a given query. We define a listwise predictor $F_\theta(E_S, R) \to \hat{\mathbf{\Phi}} \in \mathbb{R}^n$ (e.g., parameterized via a self-attention mechanism or a Transformer encoder) that outputs the predicted attribution vector for the entire context simultaneously. The ground-truth target is the exact Shapley vector $\mathbf{\Phi} = [\phi_1, \dots, \phi_n]^\top$.

\subsection{Stochastic Listwise Loss via Unbiased Estimators}

In a standard learning-to-rank setup, the ideal listwise objective is to minimize the squared Euclidean distance between the predicted vector and the true Shapley vector, $\mathcal{L}_{ideal}(\theta) = \frac{1}{2} \| \hat{\mathbf{\Phi}} - \mathbf{\Phi} \|_2^2$. However, computing the true target vector $\mathbf{\Phi}$ for every query in the training dataset reintroduces the $\mathcal{O}(2^n)$ combinatorial bottleneck. 

To resolve this, we leverage the property that the Shapley value is the expectation of marginal contributions over all permutations of the set. For a given query and its context $S$, let $\pi \sim \Pi(S)$ be a uniformly sampled permutation of the snippets. Let $S_i^\pi$ denote the subset of snippets that precede $s_i$ in the permutation $\pi$. We define the stochastic marginal gain vector $\mathbf{m}^\pi \in \mathbb{R}^n$, where each element is the closed-form geometric marginal gain of a snippet given its predecessors:
$$ m_i^\pi = \Delta v (S_i^\pi, s_i) $$
Crucially, $\mathbf{m}^\pi$ is an unbiased estimator of the true Shapley vector: $\mathbb{E}_{\pi}[\mathbf{m}^\pi] = \mathbf{\Phi}$ as it is shown in the next lemma. 


\begin{lemma}\label{lemma:unbiased_estimator}
\textbf{(Unbiased Estimator of the Shapley Vector).} 
Let $\mathbf{\Phi} = [\phi_1, \dots, \phi_n]^\top$ be the exact Shapley attribution vector for a retrieved context set $S$ evaluated under any valid characteristic function $v$ (e.g., $v_{emb}$, $v_{\log}$, or $v_{ent}$). Let $\pi$ be a permutation uniformly sampled from the symmetric group of all permutations $\Pi(S)$. Define the stochastic marginal gain vector $\mathbf{m}^\pi = [m_1^\pi, \dots, m_n^\pi]^\top$, where $m_i^\pi = \Delta v(S_i^\pi, s_i)$ and $S_i^\pi$ is the set of snippets preceding $s_i$ in $\pi$. Then $\mathbf{m}^\pi$ is a strictly unbiased estimator of the true Shapley vector, such that:
$$ \mathbb{E}_{\pi}[\mathbf{m}^\pi] = \mathbf{\Phi} $$
\end{lemma}

\begin{proofof}{Lemma \ref{lemma:unbiased_estimator}}
To prove that the vector expectation holds, we evaluate the expectation of an arbitrary $i$-th component, $\mathbb{E}_{\pi}[m_i^\pi]$. 

Because the permutation $\pi$ is sampled uniformly from the space of all possible permutations $\Pi(S)$, the probability of selecting any specific sequence is $\frac{1}{n!}$. The expected value of the marginal gain for snippet $s_i$ is therefore the arithmetic mean over all permutations:
$$ \mathbb{E}_{\pi}[m_i^\pi] = \sum_{\pi \in \Pi(S)} \frac{1}{n!} \Delta v(S_i^\pi, s_i) $$

To evaluate this sum, we group the permutations by the unordered subset of items that appear before $s_i$. Let $S' \subseteq S \setminus \{s_i\}$ be a specific subset of size $k = \vert S' \vert$. We must count how many permutations $\pi \in \Pi(S)$ satisfy exactly $S_i^\pi = S'$. 
For $s_i$ to be immediately preceded by exactly the elements in $S'$:
\begin{enumerate}
    \item The $k$ elements within $S'$ can be internally arranged in $k!$ different ways.
    \item The element $s_i$ must be fixed at position $k + 1$.
    \item The remaining $(n - k - 1)$ elements in the context can be arranged after $s_i$ in $(n - k - 1)!$ different ways.
\end{enumerate}
Therefore, there are exactly $\vert S' \vert! (n - \vert S' \vert - 1)!$ permutations that yield the exact marginal subset $S_i^\pi = S'$.

By replacing the summation over permutations with a summation over all possible subsets $S' \subseteq S \setminus \{s_i\}$ and weighting each subset by the combinatorial frequency derived above, the expectation rewrites as:
$$ \mathbb{E}_{\pi}[m_i^\pi] = \sum_{S' \subseteq S \setminus \{s_i\}} \frac{\vert S' \vert! (n - \vert S' \vert - 1)!}{n!} \Delta v(S', s_i) $$

This formulation is precisely the canonical subset-based definition of the exact Shapley value, $\phi_i$, which is agnostic to the internal mechanics of $v$. Since $\mathbb{E}_{\pi}[m_i^\pi] = \phi_i$ holds for every element $i \in \{1, \dots, n\}$, it follows that the joint vector expectation exactly recovers the full attribution vector:
$$ \mathbb{E}_{\pi}[\mathbf{m}^\pi] = \mathbf{\Phi} $$
This completes the proof.
\end{proofof}

We define the stochastic listwise loss for a single query batch as the distance between our model's joint prediction and this sampled marginal gain vector:
\begin{align} \label{eq:listwise_loss}
\mathcal{L}_{sto}(\theta) = \frac{1}{2} \left\| F_\theta(E_S, R) - \mathbf{m}^\pi \right\|_2^2 = \frac{1}{2} \sum_{i=1}^n \Big( \hat{\phi}_i - m_i^\pi \Big)^2
\end{align}



\section{Variance Reduction via Plackett-Luce Permutation Sampling (WIP), needs to be improved}

While the uniform Monte Carlo estimator detailed in Lemma \ref{lemma:unbiased_estimator} is strictly unbiased, uniform sampling across the permutation space $\Pi(S)$ can be highly inefficient. In practical RAG contexts, the vast majority of combinations contain redundant or irrelevant snippets, leading to a high variance in the sampled marginal gains and requiring a large sample size $K$ to converge. 

To achieve a tighter Mean Squared Error (MSE) for ultra-low latency inference (where $K$ must be extremely small), we can introduce a Bayesian inductive bias. By parameterizing a Plackett-Luce (PL) ranking model using an initial unbiased estimator, we can skew the sampling distribution toward permutations that are structurally more informative.

\subsection{The Plackett-Luce Sampling Distribution}

The Plackett-Luce model is a probability distribution over rankings (permutations) derived from the Luce choice axiom. Let $\hat{\mathbf{\Phi}}^{(0)} = [\hat{\phi}_1^{(0)}, \dots, \hat{\phi}_n^{(0)}]^\top$ be a fast, initial unbiased estimate of the Shapley values (e.g., computed via a small uniform sample or a pointwise geometric predictor). 

We treat these initial estimates as the latent ``skills'' or prior relevance weights of the snippets. We define a parameterized weight vector $\theta \in \mathbb{R}^n_{+}$, where:
$$ \theta_i = \exp\left( \frac{\hat{\phi}_i^{(0)}}{\tau} \right) $$
and $\tau > 0$ is a temperature hyperparameter controlling the strength of the prior.

Instead of sampling permutations uniformly ($P(\pi) = 1/n!$), we sample from the Plackett-Luce distribution $\mathcal{P}_{PL}(\pi \mid \theta)$. The probability of drawing a specific permutation $\pi = (\pi_1, \pi_2, \dots, \pi_n)$ is the sequential product of the conditional probabilities of selecting each snippet given the remaining pool:
$$ P_{PL}(\pi \mid \theta) = \prod_{j=1}^n \frac{\theta_{\pi_j}}{\sum_{k=j}^n \theta_{\pi_k}} $$

\subsection{The Biased, Low-Variance Estimator}

Using this PL distribution, we construct a new Plackett-Luce guided Shapley estimator, $\hat{\mathbf{\Phi}}^{PL}$, by taking the empirical mean of the marginal gains computed over $K$ permutations sampled strictly from $P_{PL}$:
\begin{align} \label{eq:pl_estimator}
\hat{\phi}_i^{PL} = \frac{1}{K} \sum_{k=1}^K \Delta v(S_i^{\pi_k}, s_i), \quad \pi_k \sim \mathcal{P}_{PL}(\theta)
\end{align}

By shifting the sampling distribution, we inherently violate the uniform subset combinatorial weights required by the exact Shapley formulation. Consequently, the PL-guided estimator is structurally \textbf{biased}:
$$ \text{Bias}(\hat{\phi}_i^{PL}) = \mathbb{E}_{\pi \sim PL} \big[ \Delta v(S_i^\pi, s_i) \big] - \phi_i \neq 0 $$
The estimator redefines the attribution concept from the \textit{``average marginal contribution over all possible arrivals''} to the \textit{``expected marginal contribution when highly relevant snippets are heavily biased to arrive early in the context.''}

\subsection{The Bias-Variance Tradeoff}

Despite introducing bias, the Plackett-Luce approach is theoretically highly desirable for real-time systems due to the standard Mean Squared Error (MSE) decomposition:
$$ \text{MSE}(\hat{\phi}_i^{PL}) = \text{Bias}(\hat{\phi}_i^{PL})^2 + \text{Var}(\hat{\phi}_i^{PL}) $$

In uniform sampling, the variance term dominates the MSE when $K$ is small. By concentrating the probability mass on a restricted manifold of highly probable, relevant subset trajectories, the Plackett-Luce prior acts as a powerful variance reduction mechanism. As $\tau \to 0$, the sampling distribution collapses to a deterministic ranking, driving the variance of the estimator toward zero. 

Therefore, by tuning the temperature $\tau$, the system dynamically balances the bias-variance tradeoff. For severely latency-constrained environments where $K$ is strictly bottlenecked (e.g., $K \le 5$), the bias penalty incurred by the PL estimator is mathematically eclipsed by the massive reduction in $\text{Var}(\hat{\phi}_i^{PL})$, yielding a strictly lower total MSE than the unbiased uniform estimator.

\paragraph{Optional Recovery of Unbiasedness.} 
If unbiasedness must be strictly preserved while utilizing the PL distribution to reduce variance, this formulation trivially supports classical \textbf{Importance Sampling (IS)}. By weighting each sampled marginal gain by the likelihood ratio $W(\pi) = \frac{1/n!}{P_{PL}(\pi \mid \theta)}$, the estimator recovers strict mathematical exactness ($\mathbb{E}[\hat{\phi}_i^{IS}] = \phi_i$), though at the potential risk of inflating variance if the PL prior is poorly calibrated.


%%%%%%%%%%%%%%%%%%%%%%%%%%%%%%%%%%%%%%%%%%%%%%%%%%%%%%%%%%%%%%%%%%%%%%%%



\section{Statistical Bias in Non-Linear Normalization}

The prefix-based trajectory sampling utilized in the Easton algorithm introduces a profound structural correlation among the estimated marginal gains within a single permutation batch. By the fundamental theorem of calculus applied to discrete sets, the sum of the marginal gains along any permutation trajectory perfectly telescopes to the total characteristic value of the full context:
$$ \sum_{i=1}^n m_i^\pi = v(S) - v(\emptyset) $$
Because the sum is rigidly constrained to a fixed constant $v(S)$, the individual marginal estimators $m_i^\pi$ and $m_j^\pi$ for any $i \neq j$ exhibit strict \textbf{negative covariance}. If snippet $s_i$ appears early in the permutation and claims a large marginal information gain, there is inherently less remaining semantic volume available for snippet $s_j$ later in the trajectory. 

This negative covariance interacts catastrophically with the non-linear $\max(0, \cdot)$ filter and ratio normalization if applied prematurely. Let $\hat{\phi}_i$ denote the final normalized estimator. A naive Monte Carlo implementation might apply the normalization independently to each sampled trajectory before averaging:
$$ \hat{\phi}_i^{naive} = \mathbb{E}_{\pi} \left[ \frac{\max(0, m_i^\pi)}{\sum_{j=1}^n \max(0, m_j^\pi)} \right] $$

This premature non-linear transformation introduces two distinct forms of severe statistical bias:
\begin{enumerate}
    \item \textbf{Jensen's Inequality Bias:} Because the function $f(x) = \max(0, x)$ is strictly convex, Jensen's inequality dictates that $\mathbb{E}[\max(0, X)] \ge \max(0, \mathbb{E}[X])$. If a snippet has fluctuating marginal gains (e.g., highly positive when evaluated early, but negative when evaluated late due to redundancy), the premature truncation of negative samples artificially injects phantom value into the system, artificially inflating the expectation.
    \item \textbf{Ratio Bias of Correlated Variables:} The $\max(0, \cdot)$ filter breaks the telescoping sum property. Consequently, the denominator $\sum_j \max(0, m_j^\pi)$ ceases to be a constant and becomes a random variable highly correlated with the numerator. In statistical estimation, the expectation of a ratio of correlated variables heavily diverges from the ratio of their expectations:
    $$ \mathbb{E}_{\pi} \left[ \frac{\max(0, m_i^\pi)}{\sum_j \max(0, m_j^\pi)} \right] \neq \frac{\mathbb{E}_{\pi}[\max(0, m_i^\pi)]}{\mathbb{E}_{\pi}\left[\sum_j \max(0, m_j^\pi)\right]} $$
\end{enumerate}

\begin{example}
\textbf{(Distortion from Premature Normalization).}
Consider a retrieval scenario with exactly two snippets, $S = \{s_1, s_2\}$. Assume a characteristic value function evaluating to $v(\emptyset) = 0$, $v(\{s_1\}) = 10$, $v(\{s_2\}) = 4$, and $v(\{s_1, s_2\}) = 8$. (Here, $s_2$ is individually helpful, but introduces minor contradictions when paired with $s_1$).

The exact, raw Shapley values are computed across both permutation paths:
\begin{itemize}
    \item \textbf{Path A ($s_1 \to s_2$):} $m_1^A = 10 - 0 = 10$; \quad $m_2^A = 8 - 10 = -2$.
    \item \textbf{Path B ($s_2 \to s_1$):} $m_2^B = 4 - 0 = 4$; \quad $m_1^B = 8 - 4 = 4$.
\end{itemize}
Averaging these paths yields the true raw Shapley values: $\phi_1 = \frac{10 + 4}{2} = 7$ and $\phi_2 = \frac{-2 + 4}{2} = 1$. Both are strictly positive, so the true global $\max(0, \cdot)$ filter has no effect. The \textbf{true normalized attributions} are exactly $\frac{7}{8} = 87.5\%$ for $s_1$ and $\frac{1}{8} = 12.5\%$ for $s_2$.

Now consider the biased naive Monte Carlo estimator, which applies the max filter and normalizes \textit{each path independently} before averaging:
\begin{itemize}
    \item \textbf{Path A Normalized:} The raw $[10, -2]$ becomes $[10, 0]$ after the max filter. The normalized path shares are $100\%$ for $s_1$ and $0\%$ for $s_2$.
    \item \textbf{Path B Normalized:} The raw $[4, 4]$ becomes $[4, 4]$. The normalized path shares are $50\%$ for $s_1$ and $50\%$ for $s_2$.
\end{itemize}
Averaging these normalized path shares yields an estimated attribution of $\frac{100\% + 50\%}{2} = 75\%$ for $s_1$, and $\frac{0\% + 50\%}{2} = 25\%$ for $s_2$. By aggressively truncating the $-2$ penalty in Path A before it could legitimately cancel out the $+4$ reward in Path B, the premature normalization magically injects phantom value into the system, doubling $s_2$'s apparent attribution from $12.5\%$ to $25\%$.
\end{example}

\paragraph{Asymptotically Unbiased Estimation.}
To strictly avoid these distortions, the non-linear transformations must be delayed until the empirical distribution collapses to its true mean. The theoretically sound estimator must first compute the raw empirical mean $\bar{m}_i = \frac{1}{K} \sum_{k=1}^K m_i^{\pi_k}$ over all $K$ trajectories, completely preserving the negative marginals to ensure exact expectation matching ($\mathbb{E}[\bar{m}_i] = \phi_i$). 

The normalization operators are then applied exactly once, globally, over the final aggregated means:
$$ \widehat{\text{normshap}}(i) = \frac{\max(0, \bar{m}_i)}{\sum_{j=1}^n \max(0, \bar{m}_j)} $$
By the Strong Law of Large Numbers, as $K \to \infty$, the raw empirical mean $\bar{m}_i$ converges almost surely to the exact Shapley value $\phi_i$. Because the $\max$ function and the ratio operation are continuous, the Continuous Mapping Theorem guarantees that this delayed normalization securely preserves statistical consistency, yielding an \textbf{asymptotically unbiased} estimator for the Easton attributions.


\subsection{Three-Snippet Counterexample: Distortion from Premature Normalization}

To demonstrate the bias in a three-player setting, consider a retrieval scenario with three snippets, $S = \{s_1, s_2, s_3\}$. We define the characteristic value function $v(A)$ for different subsets $A \subseteq S$ as follows:
\begin{align*}
    v(\emptyset) &= 0 \\
    v(\{s_1\}) &= 10, \quad v(\{s_2\}) = 6, \quad v(\{s_3\}) = 4 \\
    v(\{s_1, s_2\}) &= 12 \quad \text{(reflecting redundancy between $s_1$ and $s_2$)} \\
    v(\{s_1, s_3\}) &= 8 \quad \ \ \text{(reflecting contradiction/negative interaction between $s_1$ and $s_3$)} \\
    v(\{s_2, s_3\}) &= 9 \quad \ \ \text{(reflecting minor redundancy between $s_2$ and $s_3$)} \\
    v(\{s_1, s_2, s_3\}) &= 11 \quad \text{(the total characteristic value of the full context)}
\end{align*}

Because there are $3$ snippets, there are $3! = 6$ possible permutation trajectories. We calculate the raw marginal gains $m_i^\pi$ for each snippet $s_i$ along all $6$ paths:

\begin{enumerate}
    \item \textbf{Path 1 ($\pi_1: s_1 \to s_2 \to s_3$):} 
    \begin{itemize}
        \item $m_1^{\pi_1} = v(\{s_1\}) - v(\emptyset) = 10 - 0 = 10$
        \item $m_2^{\pi_1} = v(\{s_1, s_2\}) - v(\{s_1\}) = 12 - 10 = 2$
        \item $m_3^{\pi_1} = v(\{s_1, s_2, s_3\}) - v(\{s_1, s_2\}) = 11 - 12 = -1$
    \end{itemize}
    Raw gains: $[10, 2, -1]$

    \item \textbf{Path 2 ($\pi_2: s_1 \to s_3 \to s_2$):} 
    \begin{itemize}
        \item $m_1^{\pi_2} = v(\{s_1\}) - v(\emptyset) = 10 - 0 = 10$
        \item $m_3^{\pi_2} = v(\{s_1, s_3\}) - v(\{s_1\}) = 8 - 10 = -2$
        \item $m_2^{\pi_2} = v(\{s_1, s_2, s_3\}) - v(\{s_1, s_3\}) = 11 - 8 = 3$
    \end{itemize}
    Raw gains: $[10, 3, -2]$ (ordered by snippet index: $m_1=10, m_2=3, m_3=-2$)

    \item \textbf{Path 3 ($\pi_3: s_2 \to s_1 \to s_3$):} 
    \begin{itemize}
        \item $m_2^{\pi_3} = v(\{s_2\}) - v(\emptyset) = 6 - 0 = 6$
        \item $m_1^{\pi_3} = v(\{s_1, s_2\}) - v(\{s_2\}) = 12 - 6 = 6$
        \item $m_3^{\pi_3} = v(\{s_1, s_2, s_3\}) - v(\{s_1, s_2\}) = 11 - 12 = -1$
    \end{itemize}
    Raw gains: $[6, 6, -1]$

    \item \textbf{Path 4 ($\pi_4: s_2 \to s_3 \to s_1$):} 
    \begin{itemize}
        \item $m_2^{\pi_4} = v(\{s_2\}) - v(\emptyset) = 6 - 0 = 6$
        \item $m_3^{\pi_4} = v(\{s_2, s_3\}) - v(\{s_2\}) = 9 - 6 = 3$
        \item $m_1^{\pi_4} = v(\{s_1, s_2, s_3\}) - v(\{s_2, s_3\}) = 11 - 9 = 2$
    \end{itemize}
    Raw gains: $[2, 6, 3]$

    \item \textbf{Path 5 ($\pi_5: s_3 \to s_1 \to s_2$):} 
    \begin{itemize}
        \item $m_3^{\pi_5} = v(\{s_3\}) - v(\emptyset) = 4 - 0 = 4$
        \item $m_1^{\pi_5} = v(\{s_1, s_3\}) - v(\{s_3\}) = 8 - 4 = 4$
        \item $m_2^{\pi_5} = v(\{s_1, s_2, s_3\}) - v(\{s_1, s_3\}) = 11 - 8 = 3$
    \end{itemize}
    Raw gains: $[4, 3, 4]$

    \item \textbf{Path 6 ($\pi_6: s_3 \to s_2 \to s_1$):} 
    \begin{itemize}
        \item $m_3^{\pi_6} = v(\{s_3\}) - v(\emptyset) = 4 - 0 = 4$
        \item $m_2^{\pi_6} = v(\{s_2, s_3\}) - v(\{s_3\}) = 9 - 4 = 5$
        \item $m_1^{\pi_6} = v(\{s_1, s_2, s_3\}) - v(\{s_2, s_3\}) = 11 - 9 = 2$
    \end{itemize}
    Raw gains: $[2, 5, 4]$
\end{enumerate}

\subsubsection{1. True Asymptotically Unbiased Attributions}
To find the true Shapley values ($\phi_i$), we average the raw marginal gains across all $6$ paths:
\begin{align*}
    \phi_1 &= \frac{10 + 10 + 6 + 2 + 4 + 2}{6} = \frac{34}{6} \approx 5.67 \\
    \phi_2 &= \frac{2 + 3 + 6 + 6 + 3 + 5}{6} = \frac{25}{6} \approx 4.17 \\
    \phi_3 &= \frac{-1 - 2 - 1 + 3 + 4 + 4}{6} = \frac{7}{6} \approx 1.17
\end{align*}
Note that the efficiency property holds: $\sum \phi_i = 11 = v(S) - v(\emptyset)$. 

Since all true Shapley values are positive, applying the global $\max(0, \cdot)$ filter has no effect. The \textbf{true normalized attributions} are:
$$ \phi^{true} \approx [51.52\%, \ 37.88\%, \ 10.61\%] $$

\subsubsection{2. Biased Easton Estimator (Premature Normalization)}
Under Easton's implementation, the $\max(0, \cdot)$ filter and ratio normalization are applied to each path independently before averaging:
$$ \hat{\phi}_i^{Easton} = \frac{1}{6} \sum_{k=1}^6 \frac{\max(0, m_i^{\pi_k})}{\sum_{j=1}^3 \max(0, m_j^{\pi_k})} $$

We compute the normalized shares for each path:
\begin{itemize}
    \item \textbf{Path 1:} $[10, 2, -1] \to [10, 2, 0]$ (Sum = $12$) $\to [83.33\%, \ 16.67\%, \ 0\%]$
    \item \textbf{Path 2:} $[10, 3, -2] \to [10, 3, 0]$ (Sum = $13$) $\to [76.92\%, \ 23.08\%, \ 0\%]$
    \item \textbf{Path 3:} $[6, 6, -1] \to [6, 6, 0]$ (Sum = $12$) $\to [50.00\%, \ 50.00\%, \ 0\%]$
    \item \textbf{Path 4:} $[2, 6, 3] \to [2, 6, 3]$ (Sum = $11$) $\to [18.18\%, \ 54.55\%, \ 27.27\%]$
    \item \textbf{Path 5:} $[4, 3, 4] \to [4, 3, 4]$ (Sum = $11$) $\to [36.36\%, \ 27.27\%, \ 36.36\%]$
    \item \textbf{Path 6:} $[2, 5, 4] \to [2, 5, 4]$ (Sum = $11$) $\to [18.18\%, \ 45.45\%, \ 36.36\%]$
\end{itemize}

Averaging these path-level shares yields the final Easton attribution estimates:
\begin{align*}
    \hat{\phi}_1^{Easton} &= \frac{83.33\% + 76.92\% + 50.00\% + 18.18\% + 36.36\% + 18.18\%}{6} \approx 47.16\% \\
    \hat{\phi}_2^{Easton} &= \frac{16.67\% + 23.08\% + 50.00\% + 54.55\% + 27.27\% + 45.45\%}{6} \approx 36.17\% \\
    \hat{\phi}_3^{Easton} &= \frac{0\% + 0\% + 0\% + 27.27\% + 36.36\% + 36.36\%}{6} \approx 16.67\%
\end{align*}
$$ \hat{\phi}^{Easton} \approx [47.16\%, \ 36.17\%, \ 16.67\%] $$

\subsubsection{3. Analysis of Distortion}
Comparing the two results highlights the severe distortion:
\begin{table}[h]
\centering
\begin{tabular}{lccc}
\hline
\textbf{Snippet} & \textbf{True Attribution (\%)} & \textbf{Easton Attribution (\%)} & \textbf{Absolute Bias (\%)} \\ \hline
$s_1$            & 51.52                          & 47.16                            & -4.36                       \\
$s_2$            & 37.88                          & 36.17                            & -1.71                       \\
$s_3$            & 10.61                          & 16.67                            & +6.06                       \\ \hline
\end{tabular}
\end{table}

By aggressively truncating the negative marginal gains of $s_3$ in Paths 1, 2, and 3 (where it had redundant or conflicting interactions) before averaging, Easton's premature normalization prevents these penalties from offsetting the positive gains observed in Paths 4, 5, and 6. 

Consequently, $s_3$ receives a **$57\%$ relative overestimation** ($10.61\% \to 16.67\%$), demonstrating that Easton's estimator yields a statistically biased result.



\section{Transformer Attention Gram Matrices for Real-Time Attribution}

In standard Transformer-based sequence models, the self-attention mechanism computes a sequence of dense representations by dynamically routing information according to a query-key similarity matrix. Here, we formally demonstrate that the internal attention weight matrices of a trained Transformer can be directly leveraged as a Gram similarity matrix within the \textbf{Easton+} geometric framework. This formulation elegantly bridges the gap between deep sequence modeling and set-theoretic Shapley attribution, allowing real-time, zero-inference attribution of context tokens directly from a single forward pass of the model.

\subsection{Formal Formulation in Single-Head Attention}

Let $\mathbf{X} = [x_1, \dots, x_N]^\top \in \mathbb{R}^{N \times d_{\text{model}}}$ represent the input sequence matrix consisting of $N$ tokens, which contains both the retrieved context snippets and the generated response. In a single-head self-attention layer, the input sequence is projected onto three distinct vector spaces—Query ($\mathbf{Q}$), Key ($\mathbf{K}$), and Value ($\mathbf{V}$)—using learned parameter matrices $\mathbf{W}_Q, \mathbf{W}_K \in \mathbb{R}^{d_{\text{model}} \times d_k}$ and $\mathbf{W}_V \in \mathbb{R}^{d_{\text{model}} \times d_v}$:

\begin{align}
    \mathbf{Q} = \mathbf{X}\mathbf{W}_Q, \quad \mathbf{K} = \mathbf{X}\mathbf{W}_K, \quad \mathbf{V} = \mathbf{X}\mathbf{W}_V
\end{align}

The scaled dot-product attention computes the similarity score matrix before the softmax activation. Let $\mathbf{A} \in \mathbb{R}^{N \times N}$ denote the pre-softmax compatibility matrix:

\begin{align}\label{eq:raw_attention}
    \mathbf{A} = \frac{\mathbf{Q}\mathbf{K}^\top}{\sqrt{d_k}} = \frac{(\mathbf{X}\mathbf{W}_Q)(\mathbf{X}\mathbf{W}_K)^\top}{\sqrt{d_k}}
\end{align}

To deploy the embedding-based geometric framework, we must construct a valid Gram matrix $\mathbf{G} \in \mathbb{R}^{N \times N}$ directly from this attention space. While the raw pre-softmax matrix $\mathbf{A}$ is generally asymmetric ($\mathbf{A} \neq \mathbf{A}^\top$) because $\mathbf{W}_Q \neq \mathbf{W}_K$, we can define a symmetric semantic kernel by mapping the tokens into the shared query-key latent space. 

To preserve the exact bilinear forms computed by the Transformer, we define the symmetric attention Gram matrix $\mathbf{G}$ as the symmetric part of the compatibility matrix:

\begin{align}\label{eq:attention_gram}
    \mathbf{G} = \frac{1}{2} \left( \mathbf{A} + \mathbf{A}^\top \right)
\end{align}

To satisfy the geometric constraints of Lemma~\ref{lemma:vol_det} and ensure strict compliance with spatial volume metrics, we normalize the Gram matrix entries using the diagonal elements to produce a cosine-similarity equivalent kernel over the token space:

\begin{align}\label{eq:normalized_attention_gram}
    \mathbf{G}_{ij}^{\text{norm}} = \frac{\mathbf{G}_{ij}}{\sqrt{\mathbf{G}_{ii} \mathbf{G}_{jj}}}
\end{align}

By interpreting $\mathbf{G}^{\text{norm}}$ as the empirical Gram matrix of the context-response tokens, we can directly evaluate the log-volume value function $v_{\log}(S')$ for any subset of context tokens $S'$ relative to a target response token $r$ (e.g., the generated end-of-sequence or answer token). 

Because the Gram matrix $\mathbf{G}^{\text{norm}}$ is constructed directly from the attention matrices, calculating $v_{\log}(S')$ and its corresponding exact Shapley value requires zero additional neural network evaluations. The combinatorial evaluation is performed entirely as post-processing over the pre-computed attention matrix $\mathbf{G}^{\text{norm}}$ from the Transformer's existing forward pass, reducing the computational overhead of exact attribution to standard CPU-bound matrix operations. (This is not true.)

\subsection{Connection to Spectral Entropy}

When utilizing this attention-derived Gram matrix, the spectral entropy $\mathcal{H}(\rho)$ naturally quantifies the representational diversity of the sequence's attention profile. Defining the sequence density matrix as:

\begin{align}
    \rho = \frac{\mathbf{G}^{\text{norm}}}{\text{Tr}(\mathbf{G}^{\text{norm}})} = \frac{1}{N} \mathbf{G}^{\text{norm}}
\end{align}

the von Neumann entropy $\mathcal{H}(\rho) = -\sum_i p_i \log p_i$ acts as a structural measure of the sequence's attentional focus. If the attention is evenly distributed across independent, non-redundant tokens, $\mathbf{G}^{\text{norm}} \to \mathbf{I}_N$, and the spectral entropy is strictly maximized at $\log N$. Conversely, if the attention is highly redundant or concentrated on a single collinear concept, the eigenvalues of the attention Gram matrix collapse, driving the spectral entropy toward $0$.

\begin{remark}
\textbf{(Extension to Multi-Head Attention).}
In standard Transformer models, self-attention is split across $h$ independent heads, where each head $j \in \{1, \dots, h\}$ projects the sequence using distinct weight matrices $\mathbf{W}_Q^{(j)}, \mathbf{W}_K^{(j)}$ to produce head-specific pre-softmax attention matrices $\mathbf{A}^{(j)}$. To apply the Easton+ framework to a multi-head architecture, the individual attention matrices can be aggregated into a single global sequence Gram matrix through two main paradigms:
\begin{enumerate}
    \item \textbf{Linear Superposition:} A unified Gram matrix can be computed by taking a weighted average of the head-specific symmetric Gram matrices:
    \begin{align}
        \mathbf{G}_{\text{global}} = \sum_{j=1}^h w_j \mathbf{G}^{(j)}
    \end{align}
    where the weights $w_j$ can be set uniformly ($w_j = \frac{1}{h}$) or learned dynamically to reflect the relative importance of individual heads.
    \item \textbf{Concatenated Latent Representation:} Alternatively, we can construct a unified multi-head embedding space by concatenating the query and key projections across all heads for each token:
    \begin{align}
        \mathbf{z}_i^{\text{global}} = \left[ q_i^{(1)\top}, k_i^{(1)\top}, \dots, q_i^{(h)\top}, k_i^{(h)\top} \right] ^\top \in \mathbb{R}^{2h \cdot d_k}
    \end{align}
    The global Gram matrix is then directly computed as the inner product matrix of these concatenated multi-head representations: $\mathbf{G}_{\text{global}} = \mathbf{Z}^{\text{global}} (\mathbf{Z}^{\text{global}})^\top$, which naturally captures the cross-head semantic relationships of the tokens.
\end{enumerate}
\end{remark}

\section{Snippet-Level Attribution from Attention Gram Matrices}

When attributing model generation to retrieved context snippets via the attention Gram matrix $\mathbf{G} \in \mathbb{R}^{N \times N}$, a fundamental architectural decision arises: whether to compute Shapley values at the token level and aggregate post hoc, or to formulate an aggregated snippet-level characteristic function directly.

\subsection{Token-Level Aggregation vs. Snippet-Level Games}

Let $S = \{s_1, \dots, s_n\}$ denote the set of $n$ retrieved snippets, and let $\mathcal{T}_i \subset \{1, \dots, N\}$ denote the disjoint set of token indices corresponding to snippet $s_i$. Let $\mathcal{T}_r$ represent the response token indices.

\paragraph{Approach 1: Token-Level Shapley Summation.} 
One could treat all $N$ tokens as independent players in the grand game $\mathcal{N} = \{1, \dots, N\} \setminus \mathcal{T}_r$. For each token $t \in \mathcal{T}_i$, its token Shapley value is defined as:
\begin{align}
    \phi_{\text{token}}(t) = \sum_{U \subseteq \mathcal{N} \setminus \{t\}} \frac{|U|!(|\mathcal{N}| - |U| - 1)!}{|\mathcal{N}|!} \Big( v_{\log}(U \cup \{t\}) - v_{\log}(U) \Big)
\end{align}
The attribution of snippet $s_i$ is then computed via linear summation: $\Phi_{\text{sum}}(s_i) = \sum_{t \in \mathcal{T}_i} \phi_{\text{token}}(t)$.

\paragraph{Approach 2: Snippet-Level Shapley Game.} 
Conversely, we define a cooperative game directly over the $n$ snippets as atomic players, using an aggregated characteristic function $\bar{v}(S')$ for any subset $S' \subseteq S$:
\begin{align}
    \Phi_{\text{snip}}(s_i) = \sum_{S' \subseteq S \setminus \{s_i\}} \frac{|S'|!(n - |S'| - 1)!}{n!} \Big( \bar{v}(S' \cup \{s_i\}) - \bar{v}(S') \Big)
\end{align}

\begin{proposition}\label{prop:non_equivalence}
\textbf{(Non-Equivalence of Token Aggregation and Snippet Games).}
Because the geometric volume function $\text{Vol}(\cdot)$ and log-volume function $v_{\log}(\cdot)$ are strictly submodular, the summed token Shapley values do not equal the snippet-level Shapley values:
\begin{align}
    \Phi_{\text{sum}}(s_i) \neq \Phi_{\text{snip}}(s_i) \quad \text{in general.}
\end{align}
\end{proposition}

\begin{proof}
Because tokens belonging to the same snippet $\mathcal{T}_i$ exhibit strong semantic collinearity within the attention manifold, they act as mutual substitutes. In the unconstrained permutation space of $N!$ token arrivals, individual tokens from $\mathcal{T}_i$ cannibalize each other's marginal contributions: the token appearing earliest in a permutation claims the subspace expansion, while subsequent tokens from the same snippet receive vanishing marginal gains. In contrast, in the snippet game, all tokens in $\mathcal{T}_i$ arrive simultaneously as a block coalition. By the classical theory of games with coalition structures, the snippet-level Shapley value corresponds to the Owen value, which restricts the permutation distribution to block arrivals and strictly diverges from the unconstrained atomic Shapley sum whenever the underlying characteristic function possesses non-zero cross-player synergies or redundancies.
\end{proof}

Consequently, evaluating the characteristic function directly at the snippet level ($\Phi_{\text{snip}}$) is theoretically superior: it natively eliminates intra-snippet collinearity cannibalization, provides structural invariance to snippet token length, and reduces the combinatorial evaluation space from $\mathcal{O}(2^N)$ to $\mathcal{O}(2^n)$.

\subsection{Two Formulations of Snippet-Level Attention Attribution}

We propose two explicit formulations to construct the snippet-level characteristic function $\bar{v}(S')$.

\subsubsection{Algorithm Version A: Block-Gram Subspace Attribution}
Version A preserves the complete, uncompressed token-to-token attention interactions by evaluating the volume of the multi-token parallelotope spanned by all tokens belonging to the subset of snippets.

For any subset of snippets $S' \subseteq S$, let $\mathcal{T}_{S'} = \bigcup_{s_j \in S'} \mathcal{T}_j$ be the composite set of context token indices. We extract the principal submatrix of the normalized attention Gram matrix $\mathbf{G}^{\text{norm}}$ corresponding to the union of context and response tokens:
\begin{align}
    \mathbf{G}_{S'} = \mathbf{G}^{\text{norm}}_{[\mathcal{T}_{S'} \cup \mathcal{T}_r, \, \mathcal{T}_{S'} \cup \mathcal{T}_r]}
\end{align}
The block-level characteristic function is defined via the principal determinants:
\begin{align}\label{eq:v_block}
    \bar{v}_{\text{block}}(S') = \frac{1}{2} \log \det \left( \mathbf{G}^{\text{norm}}_{[\mathcal{T}_{S'} \cup \mathcal{T}_r, \, \mathcal{T}_{S'} \cup \mathcal{T}_r]} \right) - \frac{1}{2} \log \det \left( \mathbf{G}^{\text{norm}}_{[\mathcal{T}_{S'}, \, \mathcal{T}_{S'}]} \right)
\end{align}
Using Schur complements, Equation~\eqref{eq:v_block} simplifies directly to the generalized orthogonal residual of the response token subspace relative to the context subspace:
\begin{align}
    \bar{v}_{\text{block}}(S') = \frac{1}{2} \log \det \left( \mathbf{G}_{rr} - \mathbf{G}_{r, S'} \mathbf{G}_{S', S'}^{-1} \mathbf{G}_{S', r} \right)
\end{align}
where $\mathbf{G}_{S', S'} = \mathbf{G}^{\text{norm}}_{[\mathcal{T}_{S'}, \mathcal{T}_{S'}]}$ and $\mathbf{G}_{r, S'} = \mathbf{G}^{\text{norm}}_{[\mathcal{T}_r, \mathcal{T}_{S'}]}$.

\subsubsection{Algorithm Version B: Pooled Attention Kernel Attribution}
Version B projects the variable-length token representations of each snippet into a single representative centroid, yielding an exact $n \times n$ macro Gram matrix.

For each snippet $s_i$, we compute a pooled semantic latent vector $\bar{\mathbf{z}}_i \in \mathbb{R}^{2d_k}$ by taking an attention-weighted average of its constituent token query-key states $\mathbf{z}_t = [q_t^\top, k_t^\top]^\top$:
\begin{align}
    \bar{\mathbf{z}}_i = \frac{1}{|\mathcal{T}_i|} \sum_{t \in \mathcal{T}_i} \mathbf{z}_t, \quad \bar{\mathbf{E}}_i = \frac{\bar{\mathbf{z}}_i}{\|\bar{\mathbf{z}}_i\|_2}
\end{align}
Similarly, the response tokens are pooled into a single target vector $\bar{\mathbf{R}} = \frac{\bar{\mathbf{z}}_r}{\|\bar{\mathbf{z}}_r\|_2}$. We then construct an $n \times n$ snippet Gram matrix $\bar{\mathbf{G}} \in \mathbb{R}^{n \times n}$ where:
\begin{align}
    \bar{\mathbf{G}}_{ij} = \bar{\mathbf{E}}_i^\top \bar{\mathbf{E}}_j
\end{align}
The characteristic function then reduces to the standard Easton+ orthogonal distance formula:
\begin{align}
    \bar{v}_{\text{pool}}(S') = \log d_{\perp}\left(\bar{\mathbf{R}}, \, \text{span}\big(\{\bar{\mathbf{E}}_j\}_{j \in S'}\big)\right)
\end{align}

\paragraph{Comparison of Versions A and B.} 
\begin{itemize}
    \item \textbf{Version A (High Fidelity, Multi-Subspace):} Fully captures fine-grained, cross-token interactions and preserves the exact geometric rank of multi-sentence snippets. The marginal update is computed via block Schur complements in $\mathcal{O}(|\mathcal{T}_i|^3)$ operations per subset.
    \item \textbf{Version B (Ultra-Fast, Rank-1 Updates):} Condenses each snippet into a single 1D vector. While it smooths over multi-topic snippets, it unlocks exact closed-form Sherman-Morrison rank-1 updates (Equation~\eqref{eq:marginal_vlog}), evaluating the entire power set of $n \le 20$ snippets in less than $5\text{ ms}$ on a CPU.
\end{itemize}



\bibliographystyle{alpha}
\bibliography{sample}

\appendix

\section{Proof of Lemma \ref{lemma:vol_det}}
\begin{proofof}{Lemma \ref{lemma:vol_det}}
Let $E$ be the $d \times k$ matrix whose columns are the embedding vectors $E_1, \dots, E_k$. By definition, the Gram matrix is $G = E^\top E$. 

The standard geometric definition of the $k$-dimensional volume of a parallelotope formed by $k$ vectors in $\mathbb{R}^k$ is the absolute value of the determinant of the matrix formed by those vectors. Because our vectors reside in a higher-dimensional space $\mathbb{R}^d$ (where $k \le d$), we can apply an orthogonal transformation to align the $k$-dimensional subspace spanned by $E$ with the first $k$ standard basis vectors. 

Let $Q$ be a $d \times d$ orthogonal matrix (where $Q^\top Q = I$) that performs this rigid rotation, such that $Q E = \begin{bmatrix} \tilde{E} \\ 0 \end{bmatrix}$, where $\tilde{E}$ is a $k \times k$ square matrix. Because volume is invariant under rigid rotations, the $k$-dimensional volume of the parallelotope is exactly $|\det(\tilde{E})|$.

Because the inner product is invariant under orthogonal transformations, the Gram matrix remains completely unchanged:
$$ (QE)^\top (QE) = E^\top Q^\top Q E = E^\top I E = E^\top E = G $$
Simultaneously, expanding the rotated matrix block multiplication yields:
$$ (QE)^\top (QE) = \begin{bmatrix} \tilde{E} \\ 0 \end{bmatrix}^\top \begin{bmatrix} \tilde{E} \\ 0 \end{bmatrix} = \tilde{E}^\top \tilde{E} $$
Therefore, $G = \tilde{E}^\top \tilde{E}$. Taking the determinant of both sides gives:
$$ \det(G) = \det(\tilde{E}^\top \tilde{E}) = \det(\tilde{E}^\top)\det(\tilde{E}) = \det(\tilde{E})^2 $$
Taking the square root of both sides perfectly isolates the rotation-invariant volume:
$$ \sqrt{\det(G)} = |\det(\tilde{E})| = \text{Vol}(E_{S'}) $$
\end{proofof}




\end{document}
